\section{From language suggestions to an identification certificate}
\label{exec14:llm-identification}

An analyst asks whether servicing a delinquent account will increase repayment.
An assistant may propose a graph, an adjustment set or an instrument, yet these
are different objects from an identified effect. The product must preserve the
question through each transformation: which customers, which action, at which
time, compared with what alternative, and under which observation law? A formal
engine can answer the resulting mathematical question conditional on its model.
It cannot establish the institutional truth of an assumption supplied by the
language model. This division of responsibility creates a useful role for an
assistant without treating its fluency as a certificate. \citep{exec14_ident_contract}

\subsection{Evidence from different tasks should remain distinguishable}

\emph{CLadder} provides formal causal question generation and different
verbalizations, helping separate calculation from familiar real-world
associations. Its authors report improvements from structured CausalCoT
prompting and retain the difficulty of turning language into the right formal
inputs. The inspected paper and repository are separate versioned objects; a
result reported for one dataset version is not automatically a qualification of
every released generator template. In particular, a selected repository template
places an instrumental-variable ratio inside an average-treatment-effect query.
That label needs its own structural and population assumptions.
\citep{v12_la_2312_04350,exec14_ident_packages}

\emph{CounterBench} makes another contribution: 1,000 balanced binary questions
about deterministic structural mechanisms, spanning basic, joint, nested and
conditional questions. In the inspected first version, Chen and colleagues report
GPT-4o accuracy of 52.5 percent with standard prompting, 78.8 percent with
CausalCoT and 92.0 percent with CoIn. The same table reports a decrease from
51.8 to 49.8 percent for Davinci. CoIn organizes extraction, iterative reasoning
and backtracking through demonstrations; the description of an algorithm does
not by itself make the model's generated steps a machine-checked proof. The
authors identify probabilistic mechanisms and hidden confounding as extensions.
These favorable and adverse results are historical paper observations, not
measurements of this product or of current services. The inspected v1 title is
\emph{A Benchmark for Counterfactuals Reasoning in Large Language Models}; it
should not be silently conflated with later title or release variants.
\citep{exec14_ident_counterbench}

\emph{METER} instead aligns three linguistic question families over 4,145 shared
contexts. Its first level concerns causal relations expressed in text, which the
paper calls discovery. That task differs from recovering a graph or equivalence
class from an observational law. Intervention and counterfactual options are
generated by a model and then edited and filtered by human annotators. For
GPT-4o the authors report zero-shot accuracies of 87.92, 77.96 and 67.42 percent,
with corresponding CoT accuracies of 86.46, 75.14 and 66.18 percent. Its useful
lessons concern shared context, distraction by irrelevant facts, causal reversal
and unsupported statements. Curated semantic labels do not supply a complete
formal identification oracle. Comparing its headline percentage directly with
CounterBench would confound the question, answer space, information and target.
\citep{exec14_ident_meter}

Gendron and colleagues separate graph extraction from evaluating counterfactual
mechanisms. Their experiments show that providing the correct graph need not
remove mechanism-evaluation errors. They also acknowledge using model-derived
counterfactual values as proxies for some real-world ground truth. This is a
reason to retain independent numerical or logical references when evaluating a
banking workflow. A coherent graph and a confident answer can fail at different
stages. \citep{v12_la_2410_06392}

\emph{CausalDS} moves toward an agent workflow with public scene files, separate
private truth, observation noise and identification-gated questions. It
distinguishes population-average intervention targets from counterfactual targets
and conceptual variables from their noisy measurements. Those distinctions are
directly relevant to a bank's data dictionary. However, benchmark architecture
and the semantics of every implementation branch require separate inspection.
The selected source retains unknown values for missing engines and exceptions,
while returning false with error metadata for invalid target nodes. A negative
criterion result can also survive an unavailable or failed fallback. A consumer
that reads only the Boolean can turn a malformed or unresolved request into a
claim of scientific impossibility. \citep{v12_la_2607_08093,exec14_ident_packages}

Formal proof assistance introduces a further boundary. \emph{CausalSmith}
distinguishes a mechanically checked proof from a faithful formalization of the
intended statement. A changed hypothesis, a vacuous definition or an assumed
hard step can yield an artifact that fails the original scientific purpose.
The authors' internal catalog reports 14 accepted, 59 downgraded and 71 failed
runs among 144 recorded runs; the identification and structural-model clusters
contain one acceptance among 94 runs. Reruns count separately. These are not
independent estimates of autonomous research quality: the paper acknowledges
the absence of independent human quality validation and an estimated error rate
for statement matching. A valid proof of the wrong banking target remains the
wrong deliverable. \citep{exec13_census_causalsmith}

\subsection{A valid instrument need not answer the population question}

Consider a stipulated servicing experiment. At time zero an eligible population
receives a randomized encouragement $Z$. Actual servicing $A$ occurs by day
seven, and binary repayment $Y$ is measured by day thirty. Let $A(z)$ denote
take-up under encouragement $z$, and $Y(a)$ repayment under servicing $a$.
Assume encouragement is independent of the take-up type and potential outcomes,
that it affects repayment only through servicing, and that consistency and no
interference hold. Take-up is monotone: there are no borrowers who accept only
when they are not encouraged. These assumptions define the example; an
assistant's proposed instrument does not establish them. \citep{exec14_ident_exact}

Suppose one quarter of borrowers are always-takers with $(A(0),A(1))=(1,1)$,
one quarter are never-takers with $(0,0)$, and one half are compliers with
$(0,1)$. The two candidate worlds have the same take-up and differ only in
potential repayments:
\begin{center}
\begin{tabular}{lccc}
\hline
Stratum & Probability & World A: $(Y(0),Y(1))$ & World B: $(Y(0),Y(1))$ \\
\hline
Always-taker & $1/4$ & $(0,1)$ & $(1,1)$ \\
Never-taker & $1/4$ & $(0,1)$ & $(0,0)$ \\
Complier & $1/2$ & $(0,1)$ & $(0,1)$ \\
\hline
\end{tabular}
\end{center}
Under either world the complete observed joint law of $(Z,A,Y)$ is the same.
For $Z=0$, a quarter have $(A,Y)=(1,1)$ and three quarters have $(0,0)$;
for $Z=1$, these probabilities are reversed. Consequently,
\[
 \frac{E[Y\mid Z=1]-E[Y\mid Z=0]}{E[A\mid Z=1]-E[A\mid Z=0]}
 =\frac{3/4-1/4}{3/4-1/4}=1.
\]
The population average effect is nevertheless 1 in world A and $1/2$ in world B.
Exact enumeration reproduces every cell and contrast; neither world is fitted
to bank data. \citep{exec14_ident_exact}

The reason is visible algebraically. Let $C$ denote the complier stratum.
Under the stated independence and exclusion restrictions, the first-stage
difference is $P(C)$: always-takers and never-takers do not change treatment.
For the same reason the reduced-form outcome difference is
$P(C)E[Y(1)-Y(0)\mid C]$. Dividing therefore identifies the complier effect
when $P(C)>0$. The ratio supplies no new observation of the always-takers'
untreated outcome or the never-takers' treated outcome. Here their observed
outcomes and the binary restriction permit a population-effect range
$[1/2,1]$. Varying the unobserved outcomes, or mixing compatible subtypes,
attains the endpoints and intermediate values in this stipulated class.
It would be incorrect to attach the ratio's value to the unrestricted
population target. \citep{exec14_ident_exact}

The distinction can reverse an investment conclusion. Stipulate a USD 100 value
for repayment and a USD 75 servicing cost per treated borrower. A hypothetical
policy treating everyone rather than nobody then has net value between minus
USD 25 and plus USD 25 per eligible borrower. The complier calculation alone
gives plus USD 25 and cannot settle that population policy. This comparison is
neither the value of the encouragement policy nor observed bank ROI; feasibility,
actual costs and transport beyond the stipulated population remain separate
questions. The exact example explains why the query's population must survive
every language, tool and reporting conversion. \citep{exec14_ident_exact}

\subsection{Suggestion packages and formal engines have different obligations}

The selected pinned PyWhyLLM paths illustrate a useful but limited component.
Its public initialization imports suggestion classes. The main identification
suggester delegates backdoor proposals to confounder suggestions, leaves a
frontdoor method unimplemented, and retains a commented estimand path. A
separate simple suggester prompts for IV, backdoor or frontdoor names and
extracts response tags. Protocol classes declare interfaces; the validation
suggester adds prompted critique. None of these selected paths supplies a
formal identification certificate. Its selected unit tests mock model responses
and compare extracted lists, establishing intended parsing behavior rather than
live scientific validity. \citep{exec14_ident_packages}

The proposed reuse is therefore a small optional elicitation component. Its
outputs should be candidate variables or assumptions with provenance, subject
to schema checks, variable allowlists, finite retries and review. A package name
or a method called validation should not grant scientific authority. The
root MIT notice has been inspected separately from complete runtime and
dependency rights. CLadder and CausalDS likewise need version-specific
generator, observation-model and scorer audits, with their code and data rights
kept distinct. CounterBench and METER currently support the literature argument;
adopting their released datasets requires item-level semantics and rights review.
No selected external integration is enabled by this analysis.
\citep{exec14_ident_packages,exec14_ident_contract}

The formal engine needs a stricter output contract. An identified result retains
the exact query, graph, observation law, assumptions, supported model class,
expression and checkable certificate. A not-identified result needs a compatible
counterexample or an applicable complete-method certificate for the same target
and class. Unknown covers insufficient assumptions, unsupported methods, invalid
input, absent engines, timeouts and unverified negative criterion results.
Execution status and scientific status should be distinct fields, with reasons
and fallback history. For example, failure to find a backdoor set may justify
trying another identification method; it does not by itself prove that no
observational formula exists. \citep{exec14_ident_contract}

This matters in the inspected CausalDS path. Filtering an IV candidate from a
population-ATE answer avoids one target substitution, but an empty remaining
method dictionary becomes false. When a more general fallback is absent or
fails, that result still needs its criterion scope. Invalid graph-node cases
also return false alongside error metadata. The proposed wrapper should inspect
all of those fields and preserve unknown or invalid-request semantics as
appropriate. This is a static source finding and a proposed conversion rule,
not an observed failure of an installed CausalDS runtime.
\citep{exec14_ident_packages,exec14_ident_contract}

\subsection{A public prototype that tests the actual weakness}

The smallest useful prototype combines independently enumerated mechanisms with
audited formal-engine outputs. It asks the assistant to translate a banking
question, expose its assumptions, select a supported query and explain the
result without changing its scope. Same-observed-law worlds test whether the
workflow recognizes missing information. Paraphrases and renamed variables test
whether it preserves meaning. A changed population or observation mask should
trigger a changed query, while harmless wording should not. Deliberately invalid
graphs, unsupported classes and timeouts test whether uncertainty survives.
Private truth artifacts must remain separate from assistant inputs.
\citep{exec14_ident_contract}

Report false identification, target substitution, invalid certificates,
unavailable results and the reasons for abstention separately. A high answer
accuracy on identified questions cannot conceal confident answers to
unidentified ones. Measure complete costs: formalization review, tool calls,
failed executions, proof checking and human correction as well as model tokens.
These measurements can use public and generated cases before an expensive bank
experiment. They establish whether a prototype merits institutional evaluation;
they do not manufacture assignment truth or stakeholder acceptance.
\citep{exec14_ident_contract}

The engineering route is concrete. Representative rates of business-to-formal
error, reliable detection of missing assumptions and broad faithful proof
translation remain research questions. A proof checker can reduce one class of
error while leaving those questions unresolved. The product proposal should
fund a scoped prototype with explicit failure cases and retain the unresolved
claims, rather than promise that formal syntax alone solves causal reasoning.
\citep{exec14_ident_contract}
\subsection{What the complete discovery argument supports}

Language assistance can supply a candidate definition, restriction, mechanism
or analysis. Its value depends on the information that proposal adds and the
consequences of admitting it. The chapter's distinct examples show why graph
agreement, response stability, statistical evidence, identification and
monetary value require separate assessments. Retaining a compatible model
class can be the scientifically useful result even when a single answer would
be easier to report.

A defensible proposal therefore preserves the target and clock, compares
against numerical and expert alternatives, and carries semantic failures and
uncertainty through to the decision. Source-level repairs can make a tool's
behavior more faithful to its contract; they cannot create an identifying
observation. Where the remaining alternatives imply different actions, the
next question is which feasible experiment could distinguish them at worthwhile
cost. That is the subject of the next chapter.
\label{exec14:llm-identification-end}
